%% ma: L10-B20-0001
%% lop: 10
%% chuong: 7
%% bai: 20
%% dang: 10.20.1
%% loai: TLN
%% muc: VD
%% dap_an: "2,5"
%% nguon: "(NBV)-ĐỀ SỐ 9-MỖI NGÀY 1 ĐỀ THI 2025.pdf (D09, MỖI NGÀY 1 ĐỀ - Đề số 9), Phần III Câu 1"
%% kien_thuc: "Bài 19, Bài 20 (phương trình đường thẳng, giao điểm của hai đường thẳng), tỉ số đồng dạng; tia sáng qua thấu kính hội tụ (Vật lí)"
%% trang_thai: cho-duyet
%% ly_do: "Dạng mới đề xuất 10.20.1 chờ duyệt; đề dùng thuật ngữ 'phép vị tự' (không có trong SGK Toán KNTT, lời giải ở đây chỉ dùng tọa độ và tỉ số đồng dạng) và là bài toán Vật lí (thấu kính), tương tự TLN Vật lí ở D04 mà giáo viên đã cho bỏ; gợi ý sửa: thay bằng 'tỉ số A'B'/AB' hoặc bỏ câu"
%% ghi_chu: "Đáp án gốc 2,5 đúng (sympy: B' = (-7,5;10)). Hình vẽ lại theo hình gốc theo đúng tỉ lệ (gốc không theo tỉ lệ)."
\begin{ex}
Một thấu kính hội tụ có tiêu cự $OF'=OF=5$ cm (kính lúp). Vật sáng $AB=4$ cm được đặt vuông góc với trục chính của thấu kính, cách thấu kính một đoạn $OA=3$ cm, qua thấu kính cho ảnh ảo $A'B'$, trong đó $A'B'$ là ảnh của $AB$ qua phép vị tự tâm $O$ tỉ số $k$. Tính tỉ số $k$.
\begin{center}
\begin{tikzpicture}[x=.55cm,y=.4cm]
  \draw[truc] (-8.8,0)--(7,0);
  \draw[<->] (0,-3)--(0,11);
  \draw[->,thick] (-3,0)--(-3,4);
  \draw[->,thick,dashed] (-7.5,0)--(-7.5,10);
  \draw[->] (-3,4)--(0,4);
  \draw[->] (0,4)--(6.5,-1.2);
  \draw[->] (-3,4)--(0,0);
  \draw[->] (0,0)--(2.2,-2.93);
  \draw[dashed] (0,4)--(-7.5,10);
  \draw[dashed] (0,0)--(-7.5,10);
  \fill (-5,0) circle (1.3pt) (5,0) circle (1.3pt) (0,4) circle (1.3pt);
  \node[below] at (-5,0) {$F$};
  \node[above right] at (5,0) {$F'$};
  \node[below left] at (0,0) {$O$};
  \node[below] at (-3,0) {$A$};
  \node[left] at (-3,4) {$B$};
  \node[below] at (-7.5,0) {$A'$};
  \node[left] at (-7.5,10) {$B'$};
  \node[above right] at (0,4) {$H$};
\end{tikzpicture}
\end{center}
\shortans{2,5}
\loigiai{Chọn hệ tọa độ $Oxy$ có gốc $O$, trục $Ox$ là trục chính, thấu kính nằm trên $Oy$. Khi đó $A(-3;0)$, $B(-3;4)$, $F'(5;0)$.
Tia từ $B$ song song với trục chính cắt thấu kính tại $H(0;4)$, tia ló đi qua $F'$ nên nằm trên đường thẳng $HF'$: $4x+5y-20=0$. Tia từ $B$ qua quang tâm $O$ truyền thẳng, nằm trên đường thẳng $BO$: $4x+3y=0$.
Ảnh ảo $B'$ là giao điểm của hai đường thẳng này (kéo dài về phía vật): giải hệ được $B'(-7{,}5;10)$, suy ra $A'(-7{,}5;0)$.
Vì $\overrightarrow{OA'}=k\overrightarrow{OA}$ nên $k=\dfrac{-7{,}5}{-3}=2{,}5$ (kiểm tra: $A'B'=10=2{,}5\cdot4$).}
\end{ex}
